\subsection{Marked fixed-source cap-data loops}
\begin{statement}{P3K.41}{Proved}

Fix a marked rational elliptic source in the whole-zero selected
\(I_1\) sphere catalogue, its actual source line, translation
\(Q=bG\), logarithmic sites, root orders \(m_i\), norm integers
\(n_i\), and an index-one periodic fan. Retain every finite line
character and every annular map. Consider loops of the actual affine
source maps and line units of P3F.44/P3F.49 with their prescribed
linear source/lattice marking, together with the scalar \(\lambda\).
Put \(M=\prod_i m_i\), including the order-one and framing-chart
norms in
\[
           \sum_i\frac{n_i}{m_i}=\frac{\varepsilon}{M},
                \qquad\varepsilon\in\{1,-1\}.
\]
The global group of extending fibrewise translations is exactly
the physical \(\mathbb C^*\) action. In particular every such
fixed-scalar cap-data loop has trivial sphere mapping class.

More generally, a closed loop in this marked scope with scalar
winding \(\ell\) satisfies
\[
                       \ell=N\frac{M}{b},\qquad N\in\mathbb Z.
\]
Its physical return is isotopic to the \(N\)-th power of the
actual shorter scalar return of P3K.35, whose parameter is
\(\lambda=w^{M/b}\). Thus the local finite kernels of \(1-D_i\)
produce no additional return class in this specified cap-data scope.
For the transported fan coordinate \(S=p^2-pq+q^2\), with
\(\gcd(p,q)=1\), write \(H=Ma^2/d\) and let \(\chi\) be the
actual page value in P3K.35. The oriented return is
\[
 [\Sigma_\Phi]
 =\frac{N}{b}\left(\chi(4-3S)+\frac{H(3S^2-7)}8\right)
                         \pmod{28}.
\]
This equality is the integer return of a constructed shorter family
and its iterates. The proof retains the inverse closing arrow in the
physical monodromy. It concerns the stated affine cap/line/annular
variations; moving logarithmic sites, source-moduli loops, arbitrary
nontranslation endpoint maps, and the component-supported OS49
model are outside this statement.

\end{statement}
\begin{proof}
\par

\medskip\noindent\textbf{CP.1. The exact local group and its Lie sheaf}

\par\nobreak

Let \(B\) be the orbifold \(\mathbf P^1\) with the actual cap orders. Its selected \(I_1\)
point is ordinary. Let \(A\) be the sheaf of fibrewise translations of the
compactified suspension that extend on its actual root and cusp charts.
All linear lattice maps are fixed. The physical delta field gives an
exact sequence of orbifold holomorphic vector bundles

\[
 0\longrightarrow\mathcal O_B\delta\longrightarrow V
       \longrightarrow L\longrightarrow0.                 \tag{CP.1}
\]

The quotient is the source translation Lie line. Here is the exact
coarse extension, including the changed additive root model. In the
minimal Weierstrass substitution

\[
 q=z^m,\qquad x=z^{2k}X,\quad y=z^{3k}Y,
\]

the source relative differential is \(z^{-k}dX/(2Y)\). Its dual source
translation vector therefore maps to \(z^k\) times the root translation
vector. The latter has deck character \(\zeta_m^k\). An invariant holomorphic
section has the form \(z^k h(q)\) times that root vector, because \(0\leq k<m\).
It is thus exactly an original source Lie section, not a different
coarse line. The allowed \(k\) values are \(\lfloor3m/4\rfloor\), \(\lfloor2m/3\rfloor\),
\(\lfloor m/4\rfloor\), \(\lfloor m/3\rfloor\), \(\lfloor m/6\rfloor\), and \(k=1\) for the quadratic I1-star
root. These are respectively the III-star, IV-star, III, IV, II, and
I1-star substitutions already proved in F41--F49. Every one has
\(0\leq k<m\). The root character and this equality are retained under the
actual line diagrams and coefficient-generic maps of F51.

On the selected I1 the source vector is the logarithmic Tate vector,
which extends as the relative dualizing vector through the actual
toric charts. The quotient line in (CP.1) has the original coarse
extension \(\mathcal O(-1)\), since the source is rational with \(\chi(\mathcal O)=1\). The
physical delta line has coarse extension \(\mathcal O\). Taking finite invariants
is exact over \(\mathbb C\). Hence the coarse direct image of \(V\) is an extension
of \(\mathcal O(-1)\) by \(\mathcal O\), and the elementary two-chart Laurent calculation gives

\[
                       H^1(B,V)=0.                         \tag{CP.2}
\]

At a reduced retained/root fibre the source component group can be
removed from a translation lift by a suspension power: the actual
source translation \(Q\) acts transitively on those components, by the
already proved sphere restriction K15. The remaining central point
lies in the additive or multiplicative identity group and has a local
logarithm. On a good root ordinary elliptic local logarithms apply.
At the selected I1 cusp the index-one component lattice gives the
unique integral toric reindexing making an extending translation
preserve the chosen component; the remaining two torus multipliers
are holomorphic units. Their logarithms give the local Lie lift.
This is the actual local argument of P5.5, which uses the toric cover
and its deck lattice, not a claim based on the dimension of an
automorphism group. In particular

\[
 0\longrightarrow H\longrightarrow V\xrightarrow{\exp}A
       \longrightarrow0                                   \tag{CP.3}
\]

is exact as a sequence of sheaves on the orbifold atlas and its
full groupoid, where \(H\) is the extending integral period sheaf. It
is not an assertion that exponentiation is surjective on invariant
sections of a coarse root disc. A fixed torsion translation can lack
an invariant Lie logarithm; its non-equivariant atlas logarithm
has a finite-group lattice discrepancy, which remains in (CP.3).
Taking finite invariants was used only for the complex vector sheaf
\(V\) in (CP.1)--(CP.2), where averaging is exact.

\par

\medskip\noindent\textbf{CP.2. Parametric conjugacies and the discrete torsor class}

\par\nobreak

Work first over a small simply connected parameter disc. F44 and F49
give actual local source conjugations, with the line lifts and their
transported annular maps. Their constructions are analytic in that
parameter after shrinking: a central isogeny has a local inverse at
each chosen inverse image; the multiplicative square root has a local
analytic branch; the residual zero-central logarithm is solved by the
displayed finite averaging sum. The residual line unit has its
central finite character retained and its zero-central logarithm
solved by the same sum. The component choices and this finite
character are locally constant. No choice has yet been made around
the entire parameter circle.

At the cusp the scalar and the source line are fixed, so the same
toric translation model is used. Comparing all these local models
with one fixed member gives a cocycle in \(A\) on the full Cech
groupoid of the actual orbifold atlas. This includes the finite-root
group cochains as well as ordinary overlaps. Its image under (CP.3)
is a lattice-valued total two-cocycle. Local atlas logarithms can be
chosen continuously on the parameter disc; they are not presumed
invariant. Their finite-group discrepancies and ordinary overlap
discrepancies are continuous lattice-valued functions and hence are
locally constant. Subtract the reference-parameter cocycle and use
the corresponding integral total-cochain adjustment. The result is
a \(V\)-valued total one-cocycle, including its finite-group terms.
Equation (CP.2) solves this cocycle; finite averaging is valid here
because \(V\) is a complex vector sheaf. Thus no fixed torsion point is
removed by an unjustified invariant-logarithm choice.

This solution can be taken analytic in the parameter. On the two
standard source-base charts, the solution is the Cauchy/Laurent
splitting of the negative and nonnegative powers for \(\mathcal O\) and \(\mathcal O(-1)\);
finite averaging gives its orbifold version. The integrals are over
fixed compact circles, so holomorphic parameter dependence is
preserved. A splitting of the vector extension in (CP.1) also exists,
since its extension group has coarse \(H^1(\mathcal O(1))=0\); the same Laurent
calculation constructs it. Exponentiating the resulting sections
gives the actual global family isomorphism. On every local chart it
is the specified extending translation, so it transports the
annular, finite-character, line, and cusp maps literally.

Consequently a fixed-scalar cap-data family is locally isomorphic to
the fixed compactified suspension by translations. Continuing these
actual isomorphisms around a parameter interval leaves a global
extending translation at the endpoint. The finite kernel of \(1-D\)
has not been discarded: it is included in that endpoint. The next
step computes this global endpoint group.

\par

\medskip\noindent\textbf{CP.3. Lifting an endpoint away from the selected I1}

\par\nobreak

Put \(U=\mathbf P^1\setminus\{c\}\), where \(c\) is the selected
I1. We give the actual local lift, including at every reduced root or
retained fibre. On such a root disc the construction is literally
\[
              X'=N^\times/\langle B_\lambda\rangle .
                                                               \tag{CP.3a}
\]
Here \(N\) is the actual corrected multidegree-zero root line, and
\(B_\lambda\) is an invertible line map with strict uniform norm
contraction. Thus (CP.3a) is an unramified regular covering of smooth
total spaces. The same description holds at every uncapped retained
fibre, since the only zero of the Hom section is at \(c\).

Choose a holomorphic nonzero vector \(v_O(z)\) in \(N|_O\), possible
on this disc. For an extending translation \(\tau\), the map
\(z\mapsto\tau([v_O(z)])\) is an actual holomorphic section of
\(X'\). Pulling the covering (CP.3a) back to the simply connected
root disc gives a trivial covering; lifting this one section gives
a holomorphic map \(p(z)\in N^\times\). Its projection \(T(z)\)
to the source extends through \(z=0\). It meets the smooth locus
of the reduced central fibre: \([v_O(0)]\) is a regular point of
the map to the disc, a base-preserving automorphism preserves this
property, and the quotient covering is locally biholomorphic. For
the good root this assertion is immediate. For reduced II, III,
IV or \(I_n\) the regular locus of that map is exactly the smooth locus
of the reduced fibre. Compose the lifted point with a suspension
power if necessary to put its central component in the identity
component. The required power exists because the actual Q-action
is transitive on the components. This is a definite deck map in
(CP.3a), not an identification of distinct central components by hand.

The resulting section translation \(t_T\) has an invertible line
lift \(L_T\) on the entire root surface. To see the extension,
\(t_T^*N\otimes N^{-1}\) is generically trivial. Proper direct
image gives a generator whose evaluation divisor is vertical. Its
intersection with every component is zero, since both lines have
multidegree zero. The explicit component intersection matrix has
primitive kernel the whole reduced fibre. Dividing by the
corresponding integral power of its base equation therefore leaves
a nowhere-zero section. This is the actual construction of \(L_T\)
and applies equally to the II, III, IV, \(I_2\) and retained \(I_n\) charts.
For a good fibre the same argument has its single reduced component.
The finite quotient/resolution and transported-centre diagrams are
those of F41--F49; the line being used here is their actual N.

There is exact commutation with the suspension, including its line
scalar. On a smooth generic source elliptic fibre write the
specified degree-zero line as
\[
 (\mathbf C\times\mathbf C^*)/\Lambda,\qquad
 (z,v)\sim(z+\omega,\chi(\omega)v).
\]
The constant factors of automorphy give a literal lift of a source
translation, \((z,v)\mapsto(z+t,av)\). Every other holomorphic
line lift of that translation differs by a nowhere-zero holomorphic
function on the compact elliptic curve, hence by a constant. Thus
any two such lifts commute. In particular \([L_T,B_\lambda]=1\)
on every generic fibre, and holomorphic extension gives equality
on the entire reduced/root chart. Compare \(L_T(v_O)\) with the
chosen lifted point \(p\). Their ratio is a holomorphic base unit
\(c(z)\). The actual map \(c(z)L_T\) therefore has the same value
as \(\tau\) at the chosen origin of every regular compact torus.
Two translations with the same origin value agree. Equality on the
dense regular locus proves that \(c(z)L_T\) is an extending lift of
\(\tau\) on all of (CP.3a). This constructs the lift without
assuming a covering-lift theorem for an arbitrary singular fibre.

It is exactly equivariant under the finite cap action. Pullback of
a base-preserving automorphism to the actual root cover is defined
by \((x,z)\mapsto(\tau x,z)\); it commutes with its deck action.
The two line-cover maps \(A_i L A_i^{-1}\) and \(L\) therefore
differ by \(B_\lambda^{j_i}\), where \(j_i\) is an integer constant
on the connected root disc. Iterating \(m_i\) times gives
\(B_\lambda^{m_i j_i}\). But \(A_i^{m_i}=B_\lambda^{n_i}\),
and the just-proved commutation with \(B_\lambda\) makes this
iterate the identity. Strict contraction on the covering line forces
\(j_i=0\). The lift is consequently exactly equivariant for the
actual affine action; its component and finite line character have
not been omitted.

The overlaps of these lifts differ by suspension powers and give
an integral one-cocycle on \(U\) with its finite orbifold points.
Its cohomology is zero. Explicitly the orbifold fundamental group
is the free product of the finite cyclic cap groups, and every
homomorphism from it to \(\mathbf Z\) is zero. Changing the local
lifts by the resulting integral powers of \(B_\lambda\) patches
them to one lift on \(U\). Exact finite-root equivariance descends
its source point to an original source section. The actual
root/quotient/resolution diagrams give its meromorphic line map
through the original additive fibre. On an overlap the scalar
ratio is a base function, since it is constant on the compact
smooth source elliptic fibre.

Here is the selected-cusp comparison, with both integer reindexings
written. In its actual Tate and line coordinates the two period maps
have the form
\[
 \begin{split}
 h:(Z,V)&\longmapsto(q_T Z,P_T(q)V),\\
 B_\lambda:(Z,V)&\longmapsto(Q_T(q)Z,
                              q\lambda C(q)V),
 \end{split}                                                \tag{CP.3b}
\]
where \(q_T/q,P_T,Q_T,C\) are holomorphic units. The source
sections and the line divisor sections avoid the selected node;
this gives the asserted unit Tate coordinates. In the height-one
component lattice the period shifts of (CP.3b) are exactly
\((1,0)\) and \((0,1)\). This is the whole-zero I1, index-one
calculation, and is why the component-supported OS49 case is kept
outside this note.

An extending translation on the toric covering sends its chosen
component to the component with a definite index \((r,s)\).
It induces the identity on the nearby marked lattice, hence commutes
with the component deck action. Compose its lift with
\(h^{-r}B_\lambda^{-s}\). It now preserves the chosen component
and each of its boundary strata. On the open torus times the base
its map is literally
\[
                  (Z,V,q)\longmapsto(a(q)Z,b(q)V,q).
                                                               \tag{CP.3c}
\]
Evaluate it and its inverse at \((Z,V)=(1,1)\). Since both preserve
the open central torus, \(a\) and \(b\) are holomorphic units
at zero. This proves the bounded multiplier and its inverse, not
only a bound on their valuations. Conversely (CP.3c) extends in
every toric chart. The original lift is therefore
\(h^rB_\lambda^s\) times this unit map.

Its source Tate coordinate is \(q_T^r Q_T^s a(q)\). Dividing by
the actual elliptic period \(q_T^r\) leaves the same source point
and a holomorphic unit, so the source section extends through the
smooth Neron group at the selected I1. Its line multiplier is
\(P_T^r(q\lambda C)^s b(q)\); the first factor is a unit and
its exact q-order is s. Thus no line order was lost by the source
period reindexing. The global source section has meromorphic
Weierstrass coordinates on \(\mathbf P^1\), which are rational
functions. The already proved Mordell--Weil group is exactly
\(\mathbf ZG\); hence it is \(kG\) for an integer \(k\). These actual
cusp orders are the input to (CP.6), where \(j=-s\).

\par

\medskip\noindent\textbf{CP.4. The global valuation equation retains the finite kernel}

\par\nobreak

Use K35's actual meromorphic line lift \(U_G\) covering \(+G\),
and its actual rational norm ratio
\[
                   \rho=B_0^{-1}/U_G^b.                   \tag{CP.4}
\]
At a cap \(i\) write the literal local line comparison
\(U_G=z_i^{e_i}h_iL_i\), where \(h_i\) is a nonzero unit and
\(L_i\) extends invertibly. The quotient/resolution map defining
\(L_i\) is retained. At an uncapped retained fibre use its
corresponding integer order \(e_i\) and \(m_i=1\). At the
selected \(I_1\) put \(\operatorname{ord}U_G=e_c\) in its
ordinary rational frame. K35 gives the exact rational divisor orders
\[
 \operatorname{ord}_i\rho=-be_i/m_i,
 \qquad\operatorname{ord}_c\rho=-1-be_c.
\]
Their sum is zero, so
\[
                   e_c+\sum_{i\ne c}e_i/m_i=-1/b.          \tag{CP.5}
\]
The lifted global endpoint has the form \(h(t)U_G^k\). On every
root and retained fibre its invertibility gives
\(\operatorname{ord}_i h=-ke_i/m_i\). This includes the full
finite-character condition: the right hand side must be integral.
At the selected cusp apply the actual integer reindexing of CP.3.
Its source-period power has a unit line multiplier. Composing the
endpoint with \(B_\lambda^j\) removes its line order; that map
has source translation \(-jbG\) and line zero of order \(j\).
Therefore the original endpoint has line order \(-j\), with
\(j=-s\) in CP.3, and
\[
                      \operatorname{ord}_c h+ke_c=-j.
                                                               \tag{CP.6}
\]
Taking the degree of the actual rational divisor of \(h\) and
using (CP.5) gives \(j=k/b\). Hence \(k=bj\). All finite-kernel
endpoints have now been tested against the global ordinary, root
and cusp maps; no central kernel was set to zero locally.

The functions \(h\) and \(\rho^j\) have exactly the same divisor,
by (CP.4)--(CP.6), and therefore differ by a nonzero constant \(c\).
Thus the actual global line map is
\[
                hU_G^{bj}=cB_0^{-j}.
\]
On the compactified suspension \(B_\lambda=\lambda B_0\) is its
literal deck generator. The displayed map consequently acts by
physical delta multiplication \(c\lambda^j\). Conversely every
nonzero scalar is an extending physical-delta automorphism in the
source, cap and cusp charts. The scalar action is faithful: equality
with a suspension power would make a nonzero multiple of \(bG\)
the zero section, contradicting the torsion-free Mordell--Weil group.
The global extending translation group is therefore precisely
\(\mathbb C^*\).

\par

\medskip\noindent\textbf{CP.5. The actual fixed-scalar loop endpoint}

\par\nobreak

The endpoint of CP.2 is the physical scalar just computed, with every
finite central translation tested in CP.4. Choose a logarithm of
that one nonzero number \(c_{\rm end}\). Multiplication by
\(\exp(-t\log c_{\rm end})\), \(0\leq t\leq1\), is an actual
isotopy through global delta automorphisms; compose it with the
interval trivialization. Its endpoint closes literally. It is a
smooth bundle trivialization over the parameter circle, with the
same source, cap and annular identifications and the same transported
toric fan. Thus the fixed-scalar cap-data loop has zero sphere
mapping class in this scope. The scalar path acts in the actual
Cox normal coordinates by the prescribed physical circle weights;
no additional normal framing is inserted.

For a varying scalar we retain the full endpoint and compare it
with K35's actual shorter coordinate \(w\), where
\(\lambda=w^{M/b}\). The next step determines the allowed winding
by the literal cap and cusp maps.

\par

\medskip\noindent\textbf{CP.6. Exact scalar-winding obstruction and the shorter endpoint}

\par\nobreak

The arrow in this step is fixed to K35.
Let the scalar have signed winding \(\ell\in\mathbb Z\), and use
the canonical universal parameter
\(v=\log\lambda/(2\pi iM)\). Its change is
\(\alpha=\ell/M\). The comparison in CP.2 works with the actual
scalar varying as well: on each parameter disc compare to the
canonical scalar family at that same scalar. Its endpoint is a
translation isomorphism from the canonical model at \(v\) to the one
at \(v+\alpha\). All finite characters are carried with that map.

On root \(i\) the actual ordinary-to-cap parameter cochain changes by
\(M\alpha n_i\log z_i\,\delta/(2\pi i)\). Its exponential is
exactly \(z_i^{\ell n_i}\), in the arrow of WS.14--15. Hence a
source lift \(hU_G^k\) extends as an endpoint precisely when its
ordinary root multiplier has the compensating order
\[
             m_i\operatorname{ord}_i h+k e_i=-\ell n_i.
                                                               \tag{CP.7}
\]
This is an equality of actual line orders. Reducing it modulo \(m_i\)
retains the finite commutator character; it is not a rational-rank
condition. The ordinary order-one and framing charts have exactly
this formula with \(m_i=1\) and their recorded signed norm.

At the selected I1 the two endpoint source and line charts have the
same literal scalar \(\lambda\). The integral component reindexing of
CP.3 is still required; it gives
\(\operatorname{ord}_c h+ke_c=-j\), \(j\) an integer. Taking the degree
of \(h\), using (CP.5) and the actual sphere norm equation
\(\sum_i n_i/m_i=\varepsilon/M\), gives
\[
                 j=\frac{k}{b}-\frac{\varepsilon\ell}{M}.
                                                               \tag{CP.8}
\]
Consequently
\[
                  k-bj=\frac{\varepsilon b\ell}{M}
                          \quad\hbox{is an integer}.           \tag{CP.9}
\]
Since \(b\) divides \(M\) in the proved catalogue, (CP.9) forces
\((M/b)\mid\ell\). This derives the minimal closed scalar winding
from the actual cap and cusp maps even when the intermediate cap
conjugators have nontrivial finite kernels.

Conversely K35 constructs a literal endpoint for \(\ell=M/b\),
with its rational norm root, finite characters, constant delta
normalization and simultaneous Cox product. Iterate that actual
endpoint \(q=\ell/(M/b)\) times, reversing its arrow for negative
\(q\). Comparing an arbitrary cap-data endpoint of this scalar winding
with that particular endpoint now gives a global extending
translation of the same compactified family. CP.3--4 identify that
ratio with physical delta multiplication. The path in CP.5 removes
it as an actual global isotopy. Thus the return of the given
cap-data loop is the \(q\)-th power of K35's specified shorter scalar
return. No finite-kernel contribution remains untested, and no
residue-class division is used.

For completeness, the physical return uses the inverse closing
arrow. Let \(I_t:X_t\to Y_t\) be the actual interval comparison to
the canonical scalar family. Write \(E_X:X_0\to X_1\) and
\(E_Y:Y_0\to Y_1\) for their forward endpoint identifications;
\(E_Y\) is the corresponding iterate of K35's explicit endpoint.
The comparison just proved is
\[
             I_1E_XI_0^{-1}=C E_Y,                         \tag{CP.10}
\]
where \(C\) is a physical delta automorphism of \(Y_1\). Choose
parallel transport \(P_X=I_1^{-1}P_YI_0\). The physical return
is \(\Phi_X=E_X^{-1}P_X\), so the literal comparison is
\[
 I_0\Phi_XI_0^{-1}=E_Y^{-1}C^{-1}P_Y
                    \simeq E_Y^{-1}P_Y.                  \tag{CP.11}
\]
The isotopy is the global scalar path of CP.5. The right-hand map
is exactly the \(q\)-th power of the positive shorter scalar
return, with the inverse endpoint and the same parallel-transport
direction. This proves the claimed oriented sphere mapping class
without changing a closing map into a physical return by notation.

The whole argument fixes the marked source, root orders, norm
integers, logarithmic sites and fan. Its parameter transitions have
the prescribed linear source/lattice marking and vary the affine
translations, line units and transported annular data of F44/F49.
It makes no assertion about adding arbitrary nontranslation endpoint
automorphisms, moving the source or logarithmic sites, or changing
the component-supported OS49 model.

The exact certificate is \path{code/part3/fixed_source_cap_loops.py},
with its JSON and captured output. It checks every root Lie weight,
the actual valuation identities, all thirty nonempty catalogue
branches, and the free-group closing-arrow identity (CP.11).

\end{proof}
